% --- LaTeX Letter Template ---

%% EN
\documentclass[{{font_size}}, a4paper]{article}
%% PT
%\documentclass[10pt, a4paper]{article}
%\renewcommand{\baselinestretch}{1.1}

% 2. Packages (Optional but Recommended)
% inputenc: Allows you to use UTF-8 characters directly (e.g., é, ñ, ü)
% fontenc: Improves font rendering, especially for special characters
% geometry: For custom page margins (uncomment and adjust if needed)
% hyperref: Makes links clickable in PDFs (useful for email/web addresses)
\usepackage[utf8]{inputenc}
\usepackage[T1]{fontenc}
% \usepackage[margin=1in]{geometry} % Uncomment and adjust for custom margins
%\usepackage{hyperref} % Makes links clickable in PDF
\usepackage[colorlinks = true,
            linkcolor = black,
            urlcolor  = black,
            citecolor = black,
            anchorcolor = black]{hyperref}

% --------

% --- Preamble ---
%\usepackage[utf8]{inputenc}
%\usepackage[T1]{fontenc}


% Geometry for margins
%\usepackage[left=0.8in, right=0.8in, top=0.8in, bottom=0.8in]{geometry}
\usepackage[left={{margin_left}}, right={{margin_right}}, top={{margin_top}}, bottom={{margin_bottom}}]{geometry}

% Font setup (using Roboto as a good sans-serif alternative)
\usepackage{fontspec} % Requires XeLaTeX or LuaLaTeX
\setmainfont{Roboto Light} % Ensure Roboto Light is installed system-wide
\setsansfont{Roboto}     % Ensure Roboto Regular/Bold are installed system-wide
%\setmainfont{Inter Light} % Ensure Roboto Light is installed system-wide
%\setsansfont{Inter}     % Ensure Roboto Regular/Bold are installed system-wide
%\setmainfont[
%  UprightFont = *-Light,
%  ItalicFont  = *-LightItalic
%]{Inter}
%
%\setsansfont[
%  UprightFont = *-Regular,
%  BoldFont    = *-SemiBold
%]{Inter}

% If Roboto is not installed, you can try a default system font like:
% \setmainfont{Helvetica Neue Light}
% \setsansfont{Helvetica Neue}
% Or, for basic pdflatex compatibility (remove fontspec, use T1, and):
% \usepackage{helvet}
% \renewcommand{\familydefault}{\sfdefault}

% Packages for styling and layout
\usepackage{xcolor}          % For custom colors
\usepackage{ragged2e}        % For \justifying (full justification)
\usepackage{enumitem}        % For custom list item spacing
\usepackage{fontawesome5}    % For icons (e.g., phone, email, LinkedIn, calendar, star)
\usepackage{graphicx}        % For image placeholder (though we use tikz)
\usepackage{tikz}            % For drawing lines, custom elements, and the photo placeholder
\usepackage{calc}            % For calculations (e.g., minipage widths)
\usepackage{fancyhdr}        % For header/footer customization
\usepackage{pgffor}          % For the \foreach loop used in skilllist (implicitly loaded by tikz)

\usepackage{titlesec}        % For bar below section names
\usepackage{enumitem}        % Items spacing


% Define custom colors based on the image
%\definecolor{enhancvblue}{HTML}{007bff}       % Accent blue
\definecolor{enhancvblue}{HTML}{1e90ff}     % #1e90ff % Accent blue
\definecolor{enhancvgray}{HTML}{333333}     % #333333 % Dark gray for most text/headings
\definecolor{lightgraytext}{HTML}{666666}   % #666666 % Lighter gray for descriptions/details
\definecolor{mediumgray}{gray}{0.5}         % #808080
% default lightgray                         % #bfbfbf
\definecolor{lightgrayborder}{HTML}{DDDDDD} % #dddddd % For subtle borders/placeholders

% TODO: standarize colors
%\definecolor{lightgray}{gray}{0.8}
%\definecolor{myblue}{RGB}{0,102,204}


%----------------------------------------------------------------------------------------
%	AVATAR
%----------------------------------------------------------------------------------------

%\newcommand{\circleimage}[2]{
%\begin{tikzpicture}[baseline=(current bounding box.center)]
%  \clip (0,0) circle (#1);
%  \node at (0,0) {\includegraphics[width=2#1]{#2}};
%\end{tikzpicture}%
%}

% #1 = image width
% #2 = image file
{% raw %}
\newcommand{\circleimage}[2]{
\begin{tikzpicture}[baseline=(current bounding box.center)]
  \clip (0,0) circle (#1/2);
  \node at (0,0) {\includegraphics[width=#1]{#2}};
\end{tikzpicture}%
}
{% endraw %}

%% #1 = width
%% #2 = height
%% #3 = left color
%% #4 = right color
%\newcommand{\gradientbar}[4]{%
%\begin{tikzpicture}
%  \shade[left color=#3, right color=#4]
%    (0,0) rectangle (#1,#2);
%\end{tikzpicture}%
%}


%----------------------------------------------------------------------------------------
%	COLUMN LAYOUT
%----------------------------------------------------------------------------------------

\usepackage{paracol} % Required for creating multi-column layouts that can span pages automatically

\AtBeginEnvironment{paracol}{\setlength{\parindent}{0pt}} % Paracols have default paragraph indentation and this overrides it to the template default of no indentation

\setlength\columnsep{ {{columnsep}}\textwidth} % Amount of horizontal space between the columns


\pagestyle{empty} % no page numbers, headers, footers

%\setcolumnratio{0.55,0.45}
% Widths of the two columns, specified here as a ratio summing to 1 to correspond to percentages; adjust as needed for your content
\columnratio{ {{columnratio_left}}, {{columnratio_right}}} %pyarg2


%----------------------------------------------------------------------------------------
%	SECTION TITLE
%----------------------------------------------------------------------------------------

\titleformat{\section}
  {\LARGE\bfseries\scshape\raggedright}
  {}
  {0em}
  {}
  [{\titlerule[1.2pt]}]

\titlespacing*{\section}
  {0pt}    % left
  {1.2ex}  % space before title
  {0.6ex}  % space after title + rule


%----------------------------------------------------------------------------------------
%	SECTION ENTRY
%----------------------------------------------------------------------------------------

% Command for experience and education entries
\newcommand{\resumeentry}[4]{
{% raw %}
  \par\vspace{0.5em}%
%  {\textbf\Huge{\color{enhancvgray}#1}}\par% Job Title / Degree
  {\color{enhancvgray}\bfseries\Large{#1}}\par% EN
%  {\color{enhancvgray}\Large{#1}}\par% Job Title / Degree
  {\color{enhancvblue}\bfseries{#2}}\par% Company / University
  {\footnotesize%
%  \faCalendar\ \color{lightgraytext}#3 \hfill \faMapMarkerAlt\ \color{lightgraytext}#4\par% Dates and Location
%  {\color{lightgraytext}\faCalendar}\ {\color{lightgraytext}#3} \hfill {\color{lightgraytext}\faMapMarker}\ {\color{lightgraytext}#4}\par% Dates and Location
  {\color{lightgray}\faCalendar}\ {\color{mediumgray}#3} \qquad {\color{lightgray}\faMapMarker}\ {\color{mediumgray}#4}\par% Dates and Location
  }
  \vspace{0.2em}%
{% endraw %}
%  % Custom itemize for dense bullet points
%  \begin{itemize}[label=\tiny\textbullet, leftmargin=*, itemsep=0.2ex, parsep=0ex, topsep=0.2ex, partopsep=0ex]%
%    #5% Bullet points content
%  \end{itemize}%
%  \vspace{0.5em}%
}


%----------------------------------------------------------------------------------------
%	LANGUAGES SCORES (BALLS)
%----------------------------------------------------------------------------------------

% #1 = Title
% #2 = Subtitle
% #3 = Number of blue circles (0--5)
\newcommand{\scorecircles}[3]{
\noindent
\begin{minipage}[c]{0.5\linewidth}
    {\fontsize{11pt}{13pt}\selectfont \color{black} #1}\hspace{2pt}
    {\fontsize{9pt}{11pt}\selectfont \color{mediumgray} - #2}
\end{minipage}%
\hfill
\begin{minipage}[c]{0.45\linewidth}
    \raggedleft
    \begin{tikzpicture}[baseline=(current bounding box.center)]
        \foreach \i in {0,...,4} {
            \ifnum\i<#3
                \fill[enhancvblue] (\i*0.5,0) circle (0.2);
            \else
                \fill[lightgrayborder] (\i*0.5,0) circle (0.2);
            \fi
        }
    \end{tikzpicture}
\end{minipage}%
}


%----------------------------------------------------------------------------------------
%	SKILLS (TABLE)
%----------------------------------------------------------------------------------------

\usepackage{tabularx}
\usepackage{array}

% To justify text
\usepackage{ragged2e}

% Prevent awkward line breaks in second column
%\usepackage{microtype}

%\newcolumntype{R}{>{\raggedright\arraybackslash\bfseries}l}
\newcolumntype{R}{>{\bfseries}r}
\newcolumntype{Y}{>{\raggedright\arraybackslash}X}

%% Reduce horizontal padding
%\setlength{\tabcolsep}{4pt}
%
%% Vertical spacing between rows
%\renewcommand{\arraystretch}{1.1}

\setlength{\emergencystretch}{4em}
\hyphenpenalty=500
\exhyphenpenalty=500

\usepackage[none]{hyphenat}


%----------------------------------------------------------------------------------------
%----------------------------------------------------------------------------------------
%	DOCUMENT
%----------------------------------------------------------------------------------------
%----------------------------------------------------------------------------------------

\begin{document}

%\raggedright

%----------------------------------------------------------------------------------------
%	HEADER
%----------------------------------------------------------------------------------------

\hspace{-0.2in}
\begin{minipage}[t]{0.7\linewidth} % Left part of header (Name, Title, Contact Info)
  %\noindent
  \color{enhancvgray}%
  {\Huge\textbf{ {{name}}}}\par
  \vspace{0.1em}
  {\color{enhancvblue}\textbf{ {{subheader}}}}\par
  \vspace{0.5em}
  \small%
  % Contact information in a two-column tabular layout
  \begin{tabular}[t]{@{}l@{}}
    \textcolor{enhancvblue}{\faMapMarker}\ {{location}} \\
    \textcolor{enhancvblue}{\faPhone}\ {{phone}}
  \end{tabular}
  \hspace{4em} % Small horizontal space between contact columns
  \begin{tabular}[t]{@{}l@{}}
    \textcolor{enhancvblue}{\faEnvelope}\ {{email}} \\  % \faAt
%    \faMapMarkerAlt\ New York City, New York
    \textcolor{enhancvblue}{\faLinkedin}\ \href{ {{linkedin_raw}}}{ {{linkedin_clean}}}
  \end{tabular}
\end{minipage}%
\begin{minipage}[t]{0.3\linewidth} % Right part of header for photo placeholder
  \raggedleft% Align to the right
  \vspace{-1.5em} % Adjust vertical position of photo placeholder
{% if photo %}
  \circleimage{3cm}{ {{photo}} }
{% else %}
  %% Placeholder for the photo (a gray circle with "Photo" text)
  %%\tikz \draw[lightgrayborder, fill=lightgrayborder] (0,0) circle (1.5cm) node[color=enhancvgray]{\textbf{Photo}} ;
  \tikz \draw[white, fill=white] (0,0) circle (1.5cm) node[color=enhancvgray]{} ;
{% endif %}
\end{minipage}
%\vspace{1em} % Space after the entire header
\vspace{0em} % Space after the entire header


%----------------------------------------------------------------------------------------
%	SECTIONS
%----------------------------------------------------------------------------------------

\begin{paracol}{2} % Begin two-column mode


%----------------------------------------------------------------------------------------
%	FIRST COLUMN
%----------------------------------------------------------------------------------------

\section{ {{summary_name}}}

%\noindent
%{\sloppy
%}

{{summary_contents}}

{% if jobs %}
\section{ {{jobs_name}}}

%\titlerule[0.8pt]


{% for job in jobs %}
\resumeentry{ {{job.role}} }{ {{job.company}} }{ {{job.date}} }{ {{job.location}} }
{% if job['description'] %}
{{ job.description }}
{% endif %}
{% if job['items'] %}
\begin{itemize}[leftmargin=*, noitemsep, topsep=0pt]
{% for item in job['items'] %}
    \item {{item}}
{% endfor %}
\end{itemize}
{% endif %}
{% endfor %}
{% endif %}


% PASTE [[[[SCIENTIFIC PRODUCTION]]]] HERE

\switchcolumn


%----------------------------------------------------------------------------------------
%	SECOND COLUMN
%----------------------------------------------------------------------------------------

\raggedright

{% if skills %}
\section{ {{skills_name}}}

%\begin{justify}
%\end{justify}

%{\raggedright

%\begin{tabularx}{\linewidth}{>{\raggedright\arraybackslash}R X}
{\renewcommand{\arraystretch}{1.3}
\begin{tabularx}{\linewidth}{>{\raggedleft\arraybackslash}R Y}
{% for skill in skills %}
{{skill.field}} & {{skill.stack}}\\
{% endfor %}
\end{tabularx}
}
{% endif %}



{% if education %}
\section{ {{education_name}} }

{% for edu in education %}
{% if edu['link'] %}
\resumeentry{\href{ {{edu.link}}}{ {{edu.title}}}}{ {{ edu.institution }}}{ {{edu.date}}}{ {{edu.location}}}
{% else %}
\resumeentry{ {{edu.title}}}{ {{ edu.institution }}}{ {{edu.date}}}{ {{edu.location}}}
{% endif %}
{% endfor %}
{% endif %}





{% if languages %}
\section{ {{languages_name}}}

\vspace{6pt}

\noindent

{% for lang in languages %}
\begin{minipage}[t]{0.98\linewidth}
    \scorecircles{ {{lang.lang}}}{ {{lang.proficiency}}}{ {{lang.score}}}
\end{minipage}%
\vspace{4pt}
{% endfor %}
\vspace{-4pt}  % to cancel the last vspace
{% endif %}


{% if passions %}
\section{ {{passions_name}}}

\vspace{8pt}
\begin{itemize}
{% for passion in passions %}
    \item[\textcolor{enhancvblue}{ {{passion.icon}}}] {{passion.description}}
{% endfor %}
\end{itemize}
{% endif %}


% PASTE [[[[SOFTWARE DEVELOPMENT]]]] HERE

\end{paracol} % End two-column mode

\end{document}
